% Source material: README.md ("The problem"), docs/architecture.md section 2

\section{Domain and Problem Analysis}

\subsection{The centralized status quo}

Institutional elections are run on paper, on voting machines, or through a
hosted web application backed by one organiser's database. The hosted option is
usually the cheapest, fastest and easiest for voters to reach, which is why
student senates, club committees and shareholder votes tend to use it. These
are the elections we have in mind: anywhere from a few hundred to a hundred
thousand voters, with the result expected within hours. National elections are
out of scope, since they bring legal procedures, guaranteed physical access for
everyone, and protection against coercion at home, none of which this design
provides.

\subsection{Trust concentrated in a single operator}

A hosted system asks voters to trust one organiser with four separate jobs:
keeping the voter list, holding the ballot box, counting the votes and keeping
the records. Table~\ref{tab:trust} lists what a voter has to take on trust and
what \projectname{} changes about it.

\begin{table}[htbp]
  \centering
  \caption{What a voter must trust in a centralized system, and how \projectname{} removes or reduces it.}
  \label{tab:trust}
  \begin{tabular}{@{}L{6.0cm}L{9.7cm}@{}}
    \toprule
    \textbf{What the voter has to trust} & \textbf{What \projectname{} does about it}                                                                      \\
    \midrule
    The voter list is correct
                                            & Reduced, not removed: the root and the number of registered voters are
                                              fixed on-chain before voting opens and cannot change afterwards, but
                                              who is on the list is still the registrar's word.            \\
    Nobody voted twice
                                            & Each ballot carries a nullifier, and the contract rejects repeats without learning who voted. \\
    My ballot was counted
                                            & Every ballot is a transaction on the blockchain that the voter can look up.                                \\
    The count is honest
                                            & The tally is updated by contract code, not by a person.                                               \\
    The records were not edited
                                            & The records are the blockchain itself.                                                                                       \\
    \bottomrule
  \end{tabular}
\end{table}

Each of these can be abused by an ordinary administrator, and from outside it
would be very hard to tell. Fake or deleted entries in the voter list go
unnoticed, because nobody else knows how many voters there should be. Unused
votes can be filled in at the close of polling, once the administrator knows
how many are needed. And the only evidence a voter has is a confirmation
screen, shown by the same system that could have dropped their ballot.

The ``nobody voted twice'' row relies on the proof system being sound, which in
turn depends on the trusted setup described in Section~3.

\subsection{Why paper is not the answer either}

Paper ballots do take the organiser out of the ballot box and the count, but
printing, booths, a venue, polling agents, transport and hand counting are too
expensive at this scale. Paper also cannot let a voter check their own vote.
The secret ballot deliberately separates the paper from the person, and
anything that let a voter prove how they voted would also let them prove it to
someone trying to buy their vote. Watching a paper election closely also needs
agents at every booth all day. In short, paper costs money, time and the
ability to verify your own vote, while hosted voting puts all four jobs in one
pair of hands.

\subsection{Existing systems}

Several systems already move parts of an election online, and each makes a
different trade-off. Helios~\cite{helios} runs elections from a web server and
publishes cryptographic proofs that any observer can check, although the
election itself still runs on a server controlled by the organiser. Estonia's
national internet voting system was studied in a lab reproduction and found to
be open to attacks on both the voter's computer and the servers that could
change votes or results, along with serious gaps in its architecture and
procedures~\cite{estonia}. Voatz, a mobile voting app that advertised a
blockchain back end, was found to let attackers alter, block or reveal a vote;
the weaknesses were in the app and its servers, before any vote ever reached
the blockchain~\cite{voatz}.

On Ethereum, the Open Vote Network runs a small boardroom-style election
entirely inside a smart contract, with no trusted authority needed for either
the tally or privacy~\cite{openvote}. MACI~\cite{maci} goes further and resists
bribery and collusion, but relies on a coordinator who is trusted not to reveal
how people voted. \projectname{} sits between these. Like them, it keeps the
count on a public blockchain. Unlike them, it hides who voted using
zero-knowledge membership proofs and nullifiers (Section~4), and it needs no
coordinator or tallying authority. The price is that it trusts a single
registrar for the voter list, and it does not resist coercion.

\subsection{Why a blockchain, and what it does not fix}

A blockchain is slow and expensive, so it is only worth using when two things
are true. First, parties who do not trust each other need a shared record that
nobody can edit and that none of them controls. Second, anyone who trusts none
of the organisers must still be able to check that record, which means the
data and the rules have to be public. If either condition fails, an ordinary
database with access control is the cheaper choice. An election meets both
conditions: the candidates are opponents, no single record-keeper would be
accepted by all of them, and a correct result that nobody can verify does not
settle a dispute. A private blockchain shared by the authority and the
candidates would fail the second test, because only the node operators could
check it, and a candidate who joined later would have no node at all.

There are, however, four things a blockchain does not fix.

\begin{itemize}[leftmargin=1.2em,itemsep=2pt,topsep=2pt]
  \item \textbf{Eligibility.} The blockchain can check that a proof is valid,
        but not that the person behind it is a current student, or one person
        rather than three accounts. That check has to happen off-chain, where
        the registrar verifies institute credentials before accepting a
        commitment. Publishing the voter list shows that it matches the root,
        not that every entry belongs to a real person. So \projectname{} takes
        the organiser out of the ballot box and the count, but one party still
        decides who is allowed to vote.

  \item \textbf{Ballot design.} A contract cannot be changed once it is
        deployed, so an unfair rule stays unfair, whether the problem is the
        counting method, tie-breaking, the order of candidates or the wording.
        For this reason the counting rule is kept in a separate library that is
        deployed and public before any ballot is cast.

  \item \textbf{Privacy.} Everything on a public blockchain can be read by
        anyone, forever. Storing each voter's address next to their choice
        would be worse than the database it replaces, and it could never be
        deleted. Ballot secrecy therefore has to come from cryptography built
        into the design from the start.

  \item \textbf{Coercion.} Resisting coercion is a well-known hard problem in
        electronic voting~\cite{jcj}, and here our design is weaker than the
        alternatives. The same evidence that lets a voter check their vote
        could also convince someone buying their vote, whereas hosted voting
        and paper both keep that evidence away from a buyer. In
        \projectname{}, anyone who holds the voter's secret can find their
        ballot on-chain and read the choice, or simply cast the vote
        themselves first, and this works remotely and at scale. That is the
        right trade-off when the threat is a dishonest organiser, and the wrong
        one when it is a candidate with money and a small electorate.
\end{itemize}

The design in Section~4 aims for transparency and secrecy at the same time.
Eligibility is proved in zero knowledge against a voter list stored as a single
32-byte root, voters do not need to hold any cryptocurrency, and anyone can
redo the count from public data. Our cost claim is about each operation: a
ballot should cost the same gas whether the voter list is small or large, even
though the total cost of an election still grows with the number of ballots
cast. Parameter C1 in Table~\ref{tab:eval-scale} tests this in Evaluation~2.
